\documentclass[11pt]{article}
\usepackage[margin=1in]{geometry}
\usepackage{amsmath,amssymb,amsthm}
\usepackage{physics}
\usepackage{enumitem}
\usepackage{hyperref}
\usepackage{titlesec}

\theoremstyle{definition}
\newtheorem{definition}{Definition}[section]
\newtheorem{example}{Example}[section]
\newtheorem{remark}{Remark}[section]
\newtheorem{proposition}{Proposition}[section]
\newtheorem{theorem}{Theorem}[section]

\titleformat{\section}{\large\bfseries}{Lecture 15.\thesection}{1em}{}

\title{\textbf{Lecture 15: Quantum Error Correction} \\ \large Understanding Quantum Information \& Computation}
\author{Based on the lectures by John Watrous (IBM Quantum)}
\date{}

\begin{document}
\maketitle
\tableofcontents
\newpage

%----------------------------------------------------------------------
\section{Overview}
%----------------------------------------------------------------------

Building on the stabilizer formalism (Lecture~14) and the introductory examples of quantum error correction (Lecture~13), this lecture develops the general theory of quantum error-correcting codes.  Topics include:

\begin{enumerate}
    \item The general framework: codes, error models, and the Knill--Laflamme conditions.
    \item Important code families: CSS codes, the Steane code, the 5-qubit code.
    \item Bounds on code parameters (quantum Singleton and Hamming bounds).
    \item Encoding and decoding circuits.
\end{enumerate}

%----------------------------------------------------------------------
\section{General Framework}
%----------------------------------------------------------------------

\subsection{Quantum Error-Correcting Codes}

\begin{definition}
A \textbf{quantum error-correcting code} (QECC) encoding $k$ logical qubits into $n$ physical qubits is a $2^k$-dimensional subspace $\mathcal{C}\subseteq(\mathbb{C}^2)^{\otimes n}$.
\end{definition}

An \textbf{encoding} is an isometry $V:\mathbb{C}^{2^k}\to(\mathbb{C}^2)^{\otimes n}$ whose image is $\mathcal{C}$.

\subsection{Error Models}

An error model specifies a set $\mathcal{E}=\{E_0,E_1,\dots,E_r\}$ of operators (error operators) that may act on the physical qubits.  The code $\mathcal{C}$ \textbf{corrects} $\mathcal{E}$ if there exists a recovery operation (a quantum channel) $\mathcal{R}$ such that
\[
    \mathcal{R}\!\left(\sum_j p_j E_j\rho\,E_j^\dagger\right) = \rho
    \qquad\text{for all } \rho \text{ supported on } \mathcal{C}.
\]

\subsection{The Knill--Laflamme Conditions}

\begin{theorem}[Knill--Laflamme]
A code with orthonormal basis $\{\ket{c_i}\}$ for $\mathcal{C}$ can correct the error set $\mathcal{E}=\{E_a\}$ if and only if
\[
    \bra{c_i}E_a^\dagger E_b\ket{c_j} = \alpha_{ab}\,\delta_{ij}
    \qquad\forall\;i,j,a,b,
\]
where $(\alpha_{ab})$ is a Hermitian matrix depending only on the errors, \emph{not} on the codewords.
\end{theorem}

\begin{remark}
The condition says that errors either map the code space to orthogonal subspaces (distinguishable by syndrome measurement) or act identically on the code space (no correction needed).
\end{remark}

%----------------------------------------------------------------------
\section{Code Distance and the $\llbracket n,k,d\rrbracket$ Notation}
%----------------------------------------------------------------------

\begin{definition}
The \textbf{distance} $d$ of a QECC is the minimum weight of an operator that acts non-trivially on the code space but is undetectable (maps the code space to itself non-trivially).

For a stabilizer code, $d$ is the minimum weight of an element in the normaliser $\mathcal{N}(\mathcal{S})\setminus\mathcal{S}$.
\end{definition}

A code with distance $d$:
\begin{itemize}
    \item \textbf{Detects} up to $d-1$ errors.
    \item \textbf{Corrects} up to $t=\lfloor(d-1)/2\rfloor$ errors.
\end{itemize}

%----------------------------------------------------------------------
\section{Important Code Families}
%----------------------------------------------------------------------

\subsection{CSS Codes (Calderbank--Shor--Steane)}

\begin{definition}
A \textbf{CSS code} is constructed from two classical linear codes $C_1=[n,k_1,d_1]$ and $C_2=[n,k_2,d_2]$ with $C_2\subset C_1$.  The resulting quantum code is $\llbracket n,k_1-k_2,d\rrbracket$ where $d\ge\min(d_1,d_2^\perp)$ and $d_2^\perp$ is the distance of the dual code $C_2^\perp$.
\end{definition}

\paragraph{Stabilizer generators.}
\begin{itemize}
    \item For each row $h$ of the parity-check matrix of $C_1$: a generator consisting of $Z$ operators at positions where $h$ has a $1$.
    \item For each row $h'$ of the parity-check matrix of $C_2$: a generator consisting of $X$ operators at positions where $h'$ has a $1$.
\end{itemize}

CSS codes have the advantage that $X$ and $Z$ error correction can be performed independently, simplifying the decoding.

\subsection{The Steane Code $\llbracket 7,1,3\rrbracket$}

Constructed from the classical $[7,4,3]$ Hamming code $C_1$ and its dual $C_2=C_1^\perp=[7,3,4]$ (since $C_1^\perp\subset C_1$).

Stabilizer generators (6 generators for $n-k=7-1=6$):
\[
\begin{array}{ll}
g_1 = IIIXXXX & g_4 = IIIZZZZ \\
g_2 = IXXIIXX & g_5 = IZZIIZZ \\
g_3 = XIXIXIX & g_6 = ZIZIZIZ
\end{array}
\]

Logical operators: $\bar{X}=X^{\otimes 7}$, $\bar{Z}=Z^{\otimes 7}$.

The Steane code is more efficient than the Shor code ($7$ vs.\ $9$ qubits) and has nice transversal gate properties.

\subsection{The 5-Qubit Code $\llbracket 5,1,3\rrbracket$}

The smallest possible code that corrects any single-qubit error (saturates the quantum Hamming bound for $k=1$, $t=1$).

Stabilizer generators:
\[
    g_1=XZZXI,\quad g_2=IXZZX,\quad g_3=XIXZZ,\quad g_4=ZXIXZ.
\]

This is \emph{not} a CSS code.  It is a \textbf{perfect code} (every syndrome corresponds to either no error or exactly one single-qubit Pauli error).

%----------------------------------------------------------------------
\section{Bounds on Code Parameters}
%----------------------------------------------------------------------

\subsection{Quantum Singleton Bound}

\begin{theorem}
For any $\llbracket n,k,d\rrbracket$ code:
\[
    n - k \ge 2(d-1).
\]
\end{theorem}

Codes meeting this bound with equality are called \textbf{quantum MDS codes}.

\begin{example}
For $k=1$, $d=3$: $n\ge 5$.  The 5-qubit code achieves this bound.
\end{example}

\subsection{Quantum Hamming Bound}

\begin{theorem}
For any $\llbracket n,k,d\rrbracket$ code correcting $t=\lfloor(d-1)/2\rfloor$ errors:
\[
    \sum_{j=0}^{t}\binom{n}{j}3^j \le 2^{n-k}.
\]
\end{theorem}

The factor $3^j$ accounts for the three non-identity single-qubit Pauli errors at each position.

\subsection{Quantum Gilbert--Varshamov Bound}

\begin{theorem}
There exist $\llbracket n,k,d\rrbracket$ stabilizer codes whenever
\[
    \sum_{j=0}^{d-2}\binom{n}{j}3^j < 2^{n-k}.
\]
\end{theorem}

This is an \emph{existence} bound---it guarantees good codes exist but does not construct them explicitly.

%----------------------------------------------------------------------
\section{Encoding and Decoding Circuits}
%----------------------------------------------------------------------

\subsection{Encoding}

For a stabilizer code, the encoding circuit maps the $k$ logical qubits (plus $n-k$ ancilla qubits in $\ket{0}$) into the code space.  For CSS codes, the encoding can be constructed systematically from the generator matrices of the underlying classical codes.

\subsection{Syndrome Measurement}

Each stabilizer generator $g_j$ is measured using an ancilla qubit:
\begin{enumerate}
    \item Prepare ancilla in $\ket{0}$.
    \item Apply controlled-$g_j$ (a sequence of controlled-Pauli gates).
    \item Measure the ancilla $\to$ syndrome bit $s_j\in\{0,1\}$.
\end{enumerate}
The syndrome $s=(s_1,\dots,s_{n-k})$ identifies the error.

\subsection{Decoding}

\begin{enumerate}
    \item Extract syndrome $s$.
    \item Look up the most likely error $E$ consistent with $s$ (decoding table or decoder algorithm).
    \item Apply $E^\dagger$ to correct.
    \item Optionally, reverse the encoding circuit to recover the logical qubits.
\end{enumerate}

%----------------------------------------------------------------------
\section{Summary}
%----------------------------------------------------------------------

\begin{itemize}
    \item The \textbf{Knill--Laflamme conditions} give a necessary and sufficient criterion for error correctability.
    \item \textbf{CSS codes} bridge classical and quantum coding theory; the \textbf{Steane code} $\llbracket 7,1,3\rrbracket$ is a prominent example.
    \item The \textbf{5-qubit code} $\llbracket 5,1,3\rrbracket$ is the smallest single-error-correcting code (saturates the quantum Singleton bound).
    \item The \textbf{quantum Singleton}, \textbf{Hamming}, and \textbf{Gilbert--Varshamov bounds} constrain the achievable parameters $\llbracket n,k,d\rrbracket$.
    \item Encoding, syndrome extraction, and correction can all be implemented efficiently with quantum circuits.
\end{itemize}

\end{document}
